\documentclass[11pt,a4paper]{article}
\usepackage[margin=1in]{geometry}
\usepackage{amsmath,amssymb,amsthm}
\usepackage{algorithm}
\usepackage{algpseudocode}
\usepackage{booktabs}
\usepackage{hyperref}

\theoremstyle{definition}
\newtheorem{definition}{Definition}[section]
\newtheorem{theorem}{Theorem}[section]
\newtheorem{lemma}[theorem]{Lemma}
\newtheorem{remark}{Remark}[section]

\title{\textbf{Rigorous Mathematical Specification, Direct Feature Concatenation, and Gated Information Routing in DenseNets and Highway Networks}}
\author{Team Scaibu \\ \small Email: \texttt{jaydeep.9664920749@gmail.com} \\ \small Architecture and Core Engine Team}
\date{October 2026}

\begin{document}
\maketitle

\begin{abstract}
This document presents the complete mathematical formulation and algorithmic specification of Dense Connectivity (DenseNet) and Highway Gated Networks. Whereas residual networks sum representations additively ($x + \mathcal{F}(x)$), DenseNets concatenate feature maps directly across all preceding layers ($[x_0, x_1, \dots, x_l]$), maximizing feature reuse and gradient flow across $L(L+1)/2$ direct connections. In parallel, Highway Networks modulate identity flow via parametric transform and carry gates: $y = H(x) \odot T(x) + x \odot (1 - T(x))$. We formalize both topological paradigms, provide complete algorithmic pseudo-code, derive complexity bounds, calculate a numerical example, and discuss memory optimization.
\end{abstract}

\section{Notation and Mathematical Preliminaries}

Let $x_l \in \mathbb{R}^{B \times d_l}$ denote the feature map produced at layer $l \in \{0, 1, \dots, L\}$.

\begin{itemize}
    \item $k \in \mathbb{N}_{\ge 1}$: DenseNet growth rate (number of new feature channels output by each block).
    \item $x_{\le l} \triangleq [x_0, x_1, \dots, x_l] \in \mathbb{R}^{B \times (d_0 + l \cdot k)}$: Channel concatenation across all prior layer activations.
    \item $T(x) \triangleq \sigma(x W_T^T + b_T) \in (0, 1)^{B \times d}$: Highway transform gate.
    \item $C(x) \triangleq 1 - T(x) \in (0, 1)^{B \times d}$: Highway carry gate.
    \item $H(x) \in \mathbb{R}^{B \times d}$: Non-linear highway transformation.
\end{itemize}

\begin{table}[h]
\centering
\caption{Summary of Mathematical Symbols for Dense and Highway Connections}
\begin{tabular}{lll}
\toprule
\textbf{Symbol} & \textbf{Type / Domain} & \textbf{Description} \\
\midrule
$x_l$ & $\mathbb{R}^{B \times d_l}$ & Layer $l$ activation tensor \\
$k$ & $\mathbb{N}_{\ge 1}$ & DenseNet growth rate \\
$x_{\le l}$ & $\mathbb{R}^{B \times \sum d_i}$ & Concatenated cumulative feature tensor \\
$T(x)$ & $(0, 1)^{B \times d}$ & Highway transform gate vector \\
$H(x)$ & $\mathbb{R}^{B \times d}$ & Highway non-linear branch output \\
\bottomrule
\end{tabular}
\end{table}

\section{Problem Definition and Core Concepts}

\subsection{Intuition and Motivation}
In standard ResNets, summing features $x + \mathcal{F}(x)$ can impede information flow if features have different semantic representations. DenseNet resolves this by concatenating feature maps $[x_0, \dots, x_l]$, giving every layer direct access to the raw representations of all preceding layers. Highway Networks instead introduce dynamic gating: each neuron dynamically selects whether to execute a non-linear transformation ($T(x) \approx 1$) or pass the input unchanged ($T(x) \approx 0$).

\subsection{Formal Mathematical Definition}
\begin{definition}[DenseNet and Highway Equations]
\begin{itemize}
    \item \textbf{DenseNet Concatenation}:
    \begin{equation}
    x_l = H_l\left([x_0, x_1, \dots, x_{l-1}]\right), \quad x_{\le l} = [x_{\le l-1}, x_l] \in \mathbb{R}^{B \times (d_0 + l k)} \label{eq:densenet_eq}
    \end{equation}
    \item \textbf{Highway Network Gating}:
    \begin{align}
    T(x) &= \sigma\left(x W_T^T + b_T\right) \label{eq:highway_t} \\
    y &= H(x) \odot T(x) + x \odot (1 - T(x)) \label{eq:highway_y}
    \end{align}
\end{itemize}
\end{definition}

\section{Algorithmic Specification}

The computational pipeline implemented in \texttt{NnAlgoDenseHighwayConnection} is detailed below.

\begin{algorithm}[H]
\caption{DenseNet / Highway Connection Forward Evaluation}
\begin{algorithmic}[1]
\Require Layer inputs list $\{X_0, \dots, X_{l-1}\}$, current transform $H \in \mathbb{R}^{B \times d_{\text{out}}}$
\Require Mode $\in \{\text{densenet\_concat}, \text{highway\_gated}\}$, optional gate weights $W_T, b_T$
\Ensure Output tensor $Y$
\State $B \gets \text{rows}(H)$
\If{$\text{mode} = \text{densenet\_concat}$}
    \State $D_{\text{total}} \gets \sum_{i=0}^{l-1} \text{cols}(X_i) + \text{cols}(H)$
    \State $Y \gets \text{new Matrix of shape } (B, D_{\text{total}})$
    \For{$b \gets 1 \textbf{ to } B$}
        \State $\text{col} \gets 1$
        \For{$i \gets 0 \textbf{ to } l-1$}
            \For{$d \gets 1 \textbf{ to } \text{cols}(X_i)$}
                \State $Y[b][\text{col}] \gets X_i[b][d]$
                \State $\text{col} \gets \text{col} + 1$
            \EndFor
        \EndFor
        \For{$d \gets 1 \textbf{ to } \text{cols}(H)$}
            \State $Y[b][\text{col}] \gets H[b][d]$
            \State $\text{col} \gets \text{col} + 1$
        \EndFor
    \EndFor
\Else
    \State $X_{\text{in}} \gets X_{l-1}$
    \State $T_{\text{logits}} \gets X_{\text{in}} \cdot W_T^T + b_T$
    \State $Y \gets \text{new Matrix of shape } (B, d_{\text{out}})$
    \For{$b \gets 1 \textbf{ to } B$}
        \For{$d \gets 1 \textbf{ to } d_{\text{out}}$}
            \State $t \gets 1.0 / (1.0 + \exp(-T_{\text{logits}}[b][d]))$
            \State $Y[b][d] \gets H[b][d] \cdot t + X_{\text{in}}[b][d] \cdot (1.0 - t)$
        \EndFor
    \EndFor
\EndIf
\State \Return $\{\text{output\_tensor}: Y\}$
\end{algorithmic}
\end{algorithm}

\section{Theoretical Guarantees and Direct Gradient Paths}

\begin{theorem}[DenseNet Direct Gradient Paths]
In an $L$-layer DenseNet, the total number of direct connections across layers is:
\begin{equation}
\mathcal{N}_{\text{paths}} = \frac{L(L + 1)}{2}
\end{equation}
and the backward gradient $\frac{\partial \mathcal{L}}{\partial x_0} = \sum_{l=1}^L \frac{\partial \mathcal{L}}{\partial x_l} \frac{\partial x_l}{\partial x_0}$ receives direct error signals from every subsequent layer $l$ without multiplicative attenuation.
\end{theorem}

\subsection{Limitations}
\begin{itemize}
    \item \textbf{Memory Growth}: In DenseNets, total channel count grows quadratically $\mathcal{O}(L^2 k)$ with depth, requiring transition layers ($1\times 1$ convs + pooling) to compress representations between dense blocks.
\end{itemize}

\section{Complexity Analysis}

\subsection{Time and Space Complexity}
\begin{itemize}
    \item \textbf{Time}: DenseNet concatenation takes $\mathcal{O}(B \sum d_i)$ memory copy operations. Highway gating takes $\mathcal{O}(B \cdot d_{\text{in}} \cdot d_{\text{out}})$ FLOPs.
    \item \textbf{Space}: $\Theta(B \sum d_i)$ output memory for concatenated tensor.
\end{itemize}

\section{Worked Numerical Example}

Consider $B=1$, $X_0 = [1.0, 2.0]$, $H = [3.0, 4.0]$.

\begin{enumerate}
    \item \textbf{DenseNet Concatenation}:
    \begin{equation}
    Y_{\text{dense}} = [1.0, 2.0, 3.0, 4.0] \quad (\text{shape } 1 \times 4)
    \end{equation}
    \item \textbf{Highway Gating} ($T = [0.8, 0.2]$):
    \begin{align}
    Y_1 &= 3.0(0.8) + 1.0(1 - 0.8) = 2.4 + 0.2 = 2.60 \\
    Y_2 &= 4.0(0.2) + 2.0(1 - 0.2) = 0.8 + 1.6 = 2.40
    \end{align}
    $Y_{\text{highway}} = [2.60, 2.40]$.
\end{enumerate}

\section{Numerical Considerations and Edge Cases}

\begin{itemize}
    \item \textbf{Negative Gate Bias Initialization in Highway}: Initializing $b_T = -2.0$ or $-3.0$ forces $T(x) \approx \sigma(-2) \approx 0.12$, ensuring the network starts as an identity map at step 0.
\end{itemize}

\begin{thebibliography}{9}
\bibitem{huang2017densely} G.~Huang, Z.~Liu, L.~Van Der Maaten, and K.~Q.~Weinberger, ``Densely connected convolutional networks,'' in \emph{CVPR}, pp.~4700--4708, 2017.
\bibitem{srivastava2015} R.~K.~Srivastava, K.~Greff, and J.~Schmidhuber, ``Highway networks,'' \emph{arXiv preprint arXiv:1505.00387}, 2015.
\end{thebibliography}

\end{document}
